\section{Discussion}\label{sec:discussion}

% Mathematical job:
%   - Reflect on why the scoped form is the earned form (meta-limit deferred).
%   - Situate the terminal theorem across the three pillars.
%   - State channel-versus-activation as a conceptual move, not a workaround.
%   - Place the theorem next to Foundations I and the neighboring Six Birds
%     papers (Protocol Trap for P3/P6, Promotion Criteria for scoped-promotion
%     discipline).
%
% Governing sources (for author grounding, not for citation):
%   vision.md
%   paper/orientation_memo.md
%   paper/prior_literature_context.md (A1, A3, A5)
%
% Overclaim risks:
%   - Reframing the terminal theorem as stronger than earned.
%   - Importing programmatic language from vision.md.
%   - Writing philosophy rather than mathematical reflection.

The terminal theorem is stated at scoped strength over \FATCD{}.
This section discusses what the statement earns and what it leaves
unstated, how the three parts of
Sections~\ref{sec:primitive-role-architecture},
\ref{sec:admissibility-meta-theory}, and
\ref{sec:fatcd-domain}--\ref{sec:integrated-stress} combine into one
theorem stack, and how the result sits next to earlier Six Birds
literature.

\subsection{The shape of the scoped claim}\label{subsec:discussion-shape}

Each of the three scope features --- domain restriction, existence
rather than uniqueness, residual exhaustion within a covered scope
--- is structural rather than cosmetic. An unrestricted exact-six
theorem would require a domain-free argument for which the present
paper does not claim to have the instruments. A uniqueness or
canonicity claim would require a theory of translation records
identifying distinct well-formed decompositions, which lies outside
the present admissibility regime. Scope-free residual exhaustion
would require scope-free closure of the three high-priority threats
addressed in Section~\ref{sec:upper-bound-classification}. Stating
the theorem at the scope the proof supports keeps the
endogenous-instrument meta-limit (Section~\ref{subsec:meta-limit})
separate from the current claim instead of folded into it. The posture
aligns with the broader view that scientific laws form a patchwork of
domain-local regularities rather than a single global
theory~\citep{Cartwright1999}.

\subsection{The three parts as one theorem stack}\label{subsec:discussion-pillars}

The three parts introduced in Section~\ref{subsec:three-pillars} play
complementary roles. The primitive role architecture
(Section~\ref{sec:primitive-role-architecture}) fixes the vocabulary:
six role meanings, the behavioral, verification, and
structural-categorical level pattern, and the local boundary
distinctions. This vocabulary alone does not certify any claim,
since it does not specify which role talk is admissible. The
admissibility and meta-theory regime
(Section~\ref{sec:admissibility-meta-theory}) supplies that
certification through honest bookkeeping, typed non-collapse, the
forgetting and selected-lift architecture, activation thresholds,
instrument-relative visibility, and empirical-bridge admissibility.
Admissibility alone, however, yields no terminal claim; it only
regulates what may be asserted. The exact-six theorem chain
(Sections~\ref{sec:fatcd-domain}--\ref{sec:integrated-stress})
supplies the quantification: the \FATCD{} domain, the six
lower-bound witnesses, the upper-bound classification, the
decomposition/exhaustion theorem, and the integrated stress
certification together deliver Theorem~\ref{thm:scoped-exact-six}.
Vocabulary without admissibility would be undisciplined; admissibility
without a theorem chain would be inconclusive; a theorem chain without
the first two would lack a subject.

\subsection{Channel versus activation}\label{subsec:discussion-channel-vs-activation}

The channel-versus-activation discipline of
Proposition~\ref{prop:channel-vs-activation}, enforced throughout
Sections~\ref{sec:fatcd-domain}--\ref{sec:scoped-exact-six},
separates two statements that are easily conflated. The structural
statement is that every admissible description carries six typed
closure-role slots. The activation statement is that every admissible
description realizes all six slots with active projections. The
activation statement is too strong: it fails on descriptions that
carry no route grammar and therefore have no active \Pthree{}
projection, and on descriptions whose raw material carries no
\Psix{}-aligned event/provenance/replay/check witness cluster and
therefore have no active \Psix{} projection. The structural
statement, which is what we prove, is compatible with any activation
profile. Inactive channels carry one of the statuses
\(\belowthr{}\), \(\collapsedbc{}\), \(\absentrec{}\), or
\(\outsidescope{}\), describing the way a role fails to activate.
Separating the two statements is what lets the terminal theorem hold
of every \(D \in \FATCD{}\) without requiring activation of all six
roles in every description; it also lets the lower-bound theorem of
Section~\ref{sec:lower-bounds} be stated existentially over the
domain rather than uniformly. An informal reading that elides the
distinction silently recovers the stronger activation claim, which we
do not prove.

\subsection{Instrument-relative methodology}\label{subsec:discussion-instrument-relative}

A methodological choice underlying
Section~\ref{sec:admissibility-meta-theory} deserves explicit
statement as a contribution. Instruments are treated as mathematical
probes rather than as neutral tools, in the spirit of ontic structural
realism~\citep{LadymanRoss2007} and the sheaf-theoretic treatment of
contextuality~\citep{AbramskyBrandenburger2011}. The technical form
is Proposition~\ref{prop:visibility}; the associated nonclaim is that
no result here rests on an instrument-free truth.

The motivation is reflexive. The subject matter --- finite audited
typed closure descriptions and the six typed closure roles that
project them --- is the same descriptional machinery from which any
instrument for detecting closure structure must itself be built.
Operator-presentation, quotient/descent, coherence/route,
constructive/proof-theoretic, and information/regime instruments are
not vantage points outside the domain: each is a closure structure
probing another closure structure. A claim asserted under one such
instrument while silently suppressing what another would have
revealed would privilege that instrument without bookkeeping for the
privilege --- the overread that Proposition~\ref{prop:visibility}
blocks.

The consequence is epistemic. Every admissible claim carries an
explicit visibility record: its visible content, its suppressed
content, the instruments under which visibility and suppression are
asserted, and the resulting instrument-relative claim strength
(Proposition~\ref{prop:visibility}). Triangulating a role projection
across several instruments is a disciplined way to strengthen an
assertion; reading a projection off a single instrument and declaring
its content canonical is not admissible. The discipline is one of
epistemic hygiene, in the spirit of Bohr's complementarity
doctrine~\citep{Bohr1958}: we do not claim that instruments are
incompatible in principle, only that they are not interchangeable
without explicit visibility bookkeeping and, where contexts shift,
explicit translation records
(Propositions~\ref{prop:visibility} and~\ref{prop:level-profile}).

This posture also explains why the endogenous-instrument meta-limit
(Sections~\ref{subsec:meta-limit} and~\ref{subsec:future-meta-limit})
is deferred to future work. Taking instrument-relativity seriously
from the outset lets the meta-limit be stated as its own question ---
whether the accepted internal instrument family can certify its own
soundness and completeness --- a question made statable, but not
settled, by the stance adopted here.

\subsection{Relation to earlier Six Birds results}\label{subsec:discussion-earlier-results}

Six Birds Foundations~I~\citep{Tsiokos2026SixBirdsFoundations}
introduced the theory-package formalism
\(\mathcal{T} = (Z, f, \Sigma_f, E, \mathcal{A})\) and derived the
\Pone{}--\Psix{} vocabulary under minimal hypotheses of composability
with limited interface access. The present paper does not re-derive
that vocabulary; it takes the six roles as given and asks a
different question: whether, within a specific admissibility regime
and over a specific class of admissible descriptions, the six roles
are exactly enough. The Foundations theorem gives the canonical
unavoidability side of the picture; the terminal theorem here gives
the scoped exhaustion side. The two statements address different
questions and address them at different levels; combining them does
not yield a stronger claim, and neither theorem follows from the
other.

The theorem chain also draws on a nearby
result. \citet{Tsiokos2026SixBirdsProtocolTrap}, complementing the
no-go analysis of \citet{Tsiokos2026SixBirdsNoGo}, established that
protocol holonomy is not a directionality certificate:
noncommutativity of update sequences can be nonzero while
path-reversal asymmetry remains zero, so route mismatch is a
diagnostic, not an arrow of time. That observation is load-bearing
for the role-projection alignment rules of
Section~\ref{sec:upper-bound-classification}, where \Pthree{}
route/coherence mismatch is kept structurally distinct from the
affinity-carrying face of \Psix{}. The next section records the
scope boundaries and nonclaims of the present theorem stack.
